% results-information: Q4 (where information lives), Q5 (operator levers), Q6 (thermodynamics)
% Style: ASD-STE100 Simplified Technical English (2026-10-04 rewrite).
\section{Where information lives, and what an operator can do}
\label{sec:information}

This section treats three questions.
Q4 asks where the information of the swarm lives and what it is worth.
Q5 asks what an operator can change, and Q6 asks if the swarm produces entropy beyond its parts.
All numbers come from non-holdout goal periods, and no result in this section has passed a holdout test.
The one holdout test so far, prediction C1 of H04 for the response to a change in working hours, failed with the opposite sign.

\subsection{The context window carries the value}
\label{sec:info-context}

Semantic information is the part of the system--environment mutual information whose scramble lowers viability~\cite{in:kolchinsky2018}.
Viability $V$ is the work commits of an agent per 20 calls, and $I_c$ is the information of channel $c$ about the repo of the next commit.
A natural scramble destroys $I_c$ and changes $V$ by $\Delta V_c$, and $\kappa_c=\Delta V_c/I_c$ is the value per bit (model 04;~\cite{in:sowinski2023}).
The best scramble is the forced erasure: the scaffold erases the context window at 41 records (40 calls), whatever the task state.
Thus the time of an erasure is quasi-random relative to the work (natural experiment NE41).

The erasure cost appears at three scopes.
Over the next 10 calls, a forced erasure cuts work commits by $-39\%$ $[-42,-35]$ in 7/9 regime-III goal periods (H15)\cred{0.78}.
A second estimator gives $-38\%$ $[-41,-35]$ on the same nine goal periods (H44)\cred{0.88}.
Over the next 20 calls, the cost is 41\% (H87, below).
Over a 40-call segment, the erasures cost 4.8--10.9\% of the commits in that segment (H44).
Notes that the agent writes to memory at the erasure do not make the dip shorter (within-agent $\rho=-0.03$ in G51).

After the erasure, the agent reads again in order and does not thrash (Fig.~\ref{fig:info-H44}).
In G51, the share of calls that read files, git, notes or the screen rises from 0.22 to 0.47 and relaxes over about 7 calls $[5.4, 9.2]$.
The conditional rise is $\Theta_c=+0.079$ $[0.071, 0.086]$, with the CI above 0 in 9/9 goal periods.
Category entropy falls ($-0.027$ $[-0.050,-0.004]$), so the erasure acts as a field on the call mix, not as a temperature pulse.

\begin{figure*}
\includegraphics[width=\textwidth]{H44-context-wipe/fig.pdf}
\caption{A forced erasure costs output and starts an ordered phase of reads (H15, H44).
Panels (b)--(d) align G51 calls to 14{,}593 forced erasures (vermillion) and to a boundary with no erasure (gray, the null), with 95\% agent-day bootstrap bands.
(e)~Work commits over calls $+1\ldots+10$ fall by $-38\%$ $[-41,-35]$ over nine goal periods. (f)~A command loop recurs much less often after an erasure.}
\label{fig:info-H44}
\end{figure*}

Figure~\ref{fig:info-H87} gives the $\kappa$ table over 18{,}760 forced erasures in nine goal periods (H87).
The erasure destroys $I_C=0.078$ $[0.051, 0.102]$ bits and costs $\Delta V_C=0.41$ $[0.38, 0.44]$ commits per 20 calls.
Thus $\kappa_C=5.2$ $[3.8, 7.9]$ commits per 20 calls per bit\cred{0.63}, and the first run of the estimator in H70 gave $5.2$ $[3.8, 8.4]$\cred{0.70}.
The own artifact carries the most bits ($I_A=0.14$ $[0.12, 0.16]$), but its value is zero within error ($\kappa_A=-0.6$ $[-1.6, +0.3]$).
Thus the artifact tells the agent where to work, and the context window holds what makes the work go.
Memory notes, chat reads, history search and human messages each carry at most 0.04 bits, so their $\kappa$ is not identified.

\begin{figure*}
\includegraphics[width=\textwidth]{H87-kappa-table/fig.pdf}
\caption{The $\kappa$ table: only the context window buys commits per bit (H87, with H70).
Each row is a channel at its natural scramble (NE41 forced erasures; the kickoff at NE34), with agent-day paired bootstrap 95\% intervals.
$\kappa_C=5.2$ $[3.8, 7.9]$ commits per 20 calls per bit. The own artifact carries the most bits, but $\kappa_A=-0.6$ $[-1.6,+0.3]$. The label ``n.i.'' marks rows below the 0.02-bit identification floor.}
\label{fig:info-H87}
\end{figure*}

Chat received after an erasure increases output by 22\% $[8, 36]$ without allocation bits (H87).
This is a candidate lever, but we did not test it as an intervention.
The context window also holds the loops of the swarm.
A statement in context is near-copied $3.38\times$ $[2.56, 4.46]$ more often than an erased statement, in 4/4 scorable goal periods (H69)\cred{0.85}.
A forced erasure increases the loop exit rate $\times2.40$ $[1.18, 4.88]$, so restatement loops are fixed points of the context window.

History search has no measured value.
In the search outage of 2026-03-31/04-01, commits per 20 calls changed by $+0.03\pm0.51$ (H84)\cred{0.42}.
The data exclude a dip of 0.30 or more at power 0.85--0.91, but the outage fully reached only one agent.

\subsection{Idle dynamics: traps age}
\label{sec:info-traps}

An idle agent waits at a timer gate, and each call after an idle call is a chance to act.
Kramers escape from one trap depth gives a constant escape probability, and a broad distribution of depths gives one that falls with trap age $a$.
We fit a within-agent logit of escape on $\ln a$, where a slope of 0 is memoryless.

Traps age.
The point slope is below $-0.3$ in 6/8 regime-III goal periods, and the CI is below $-0.3$ in 3/8 (H16)\cred{0.70}.
The aging continues after the nudger stops (NE43: $\beta=-0.59$ $[-0.68,-0.39]$).
Aging is not input starvation: the slope survives an input-starvation clock in 4/4 powered goal periods (H72)\cred{0.78}.
Recent input lowers escape ($\beta_s=+0.23$ $[+0.11,+0.33]$ in G51), so ambient chatter is a field that keeps agents idle.

At the gate, an unaddressed agent message decreases the odds of action to 0.50 $[0.23, 0.95]$ against no new item (H09, regime~III)\cred{0.56}.
A new @-mention gives an odds ratio of 3.09 $[1.95, 4.50]$, and 1.20 $[0.76, 1.87]$ against no new item.
Thus a mention restores the rate with no new item.

\subsection{Operator levers}
\label{sec:info-levers}

An operator input can act as a field, which changes the state of an agent, or as a catalyst, which changes only exit rates.
The auto-nudger reads the agent state and chooses a target, as a Maxwell demon does.
Its efficiency $\eta_{\rm SU}$ compares the work that it buys with the frontier $V^*(R)$, the maximum work that the same bits can buy (model 04).

A first nudge adds 1.16 $[0.75, 1.57]$ active minutes per 30 minutes to its target in G51 (H04)\cred{0.49}.
H30 finds 0.98 $[0.63, 1.37]$ in G51 and 0.95 $[0.38, 1.54]$ in G38, and about 0 for room-mates\cred{0.39}.
No lever is a pure catalyst: nudges increase idle escape $\times1.3$--$1.6$ and decrease the idle share by 0.048 $[0.019, 0.076]$ (H39)\cred{0.67}.
A second input that a later call reads keeps 80--130\% of the effect of the first input, at spacings from 1\,min to 4\,h (H43)\cred{0.69}.
Only inputs that the same call reads are wasted.

The nudger reads 1.42 bits of agent state per nudge in G51, mostly trap age (H35)\cred{0.61}.
Its bits buy glances near the frontier: $P(\text{active within 30 min})$ rises by 0.082 $[0.035, 0.132]$ per first nudge, at $\eta_{\rm SU}=0.88$ $[0.26, 0.97]$.
For committed work, it buys $+0.22$ $[-0.29, +0.82]$ commits per hour per nudge, at most about 0.4\% of swarm commits.
The switch-on (NE10) and the switch-off (NE43) of the nudger do not move swarm work.
The demon wastes its bits: 38\% of nudges go to traps of 10 or more pauses, where a nudge buys 0.5 active minutes, against 2.2--2.6.

At the logged budget, one nudge at the second idle gate of each trap gives $\times1.55$ $[1.25, 1.81]$ the extra active calls of the logged nudger (H60, G51)\cred{0.44}.
A random gate gives $\times1.46$ $[1.22, 1.67]$, and a fitted index does not beat the rule.
A nudge read at an idle gate adds 7.2 $[5.0, 9.8]$ active calls in 30 minutes.
For commits, the gain of the rule is $\times1.29$, with a CI that includes no gain (H35).
A read correction changes blocked-spell escape by $+0.007$ $[-0.053, +0.070]$ per 5-minute window over 337 windows (H55, round 1c)\cred{0.57}.
The power for a 1-logit effect is 0.96, so corrections do not free blocked agents.

\subsection{Thermodynamics: the arrow of time is individual}
\label{sec:info-thermo}

Entropy production (EP) measures how much the forward statistics of the action-class sequence of an agent differ from the time-reversed statistics~\cite{in:aguilera2026}.
A held-out Newton estimator gives a lower bound on EP, which we test against a reversible block-flip null.
The excess part of EP~\cite{in:kolchinsky2026} sets a Wasserstein speed limit $T\ge W/\bar A$ on the time to move an allocation a distance $W$ at activity $\bar A$ (model 15).

Fine-action irreversibility is real, and most of it belongs to the agents.
At least 80\% of agents exceed the reversible null in 34/35 non-holdout goal periods (H14; H56 recheck)\cred{0.82}.
The scaffold share of the excess is 0.13 in regime~III, against 0.45 in regime~I (H56).
Lab differences survive the removal of scaffold records ($\eta^2=0.42$, $p=0.0005$), and no class of scaffold change moves EP (H56)\cred{0.69}.
In behavior, no collective excess beyond the parts survives Holm correction (largest: G51, 0.009\,nats/min, Holm $p=0.20$; H14).

The first and last 30 minutes of the day carry 80--100\% of the debiased excess EP of the swarm, against a 12--25\% time share (H76)\cred{0.73}.
The runner starts the agents together, and they stop at different times, so this excess is a scheduler field.

A kickoff that names its target moves the swarm at the speed limit.
In G39, G40 and G51, the slack $S=\bar A T/W$ is 1.00--1.40, and the allocation settles in 0.25--0.5 active hours (H75)\cred{0.47}.
A free-choice goal settles with more slack: G41 takes 9.5\,h at $S=5.1$ $[1.4, 11.2]$.
On the same days in G44, the named-target room settles in 0.75\,h ($S=2.5$) and the free-choice room in 4.25\,h ($S=14.5$).

Repo dynamics show no selection beyond neutral copies (H77)\cred{0.55} and no catalytic cycle of length 3 or more (H79)\cred{0.72}~\cite{in:pinero2025,in:hordijk2023}.

The one collective arrow of time that we found goes through address.
In G51 talk, an agent starts to talk in the minute after an agent who named it ($\sigma_{\rm nam}=7.6$ $[4.3, 11.2]\times10^{-3}$\,nats/min, $p=0.01$; H90)\cred{0.41}.
Unnamed partners carry $-0.4$ $[-1.7, 0.5]\times10^{-3}$\,nats/min.
The seven other goal periods have power of at most 0.45 at this effect size, so the scope is G51 only.

\subsection{What an operator can do}
\label{sec:info-operator}

The rules below follow from the results above. No rule is holdout-confirmed.
\begin{enumerate}
\item Count each forced erasure as a cost: 39\% of commits over 10 calls (H15), 41\% over 20 calls (H87) and 4.8--10.9\% of a segment (H44). The memory file does not recover the cost. We did not measure the best call-cap length.
\item Use an erasure to stop a command loop. In G38, a loop recurs after 16\% of erasures, against 72\% with no erasure.
\item Address an idle agent by name. Unaddressed chatter halves its odds of action (odds ratio 0.50), and a mention restores them.
\item Nudge each trap one time, at its second idle gate. Do not nudge deep traps again. Expect about 7 extra active calls per read nudge, not extra commits.
\item Do not expect a correction to free a blocked agent (escape change $+0.007$, power 0.96).
\end{enumerate}
